\begin{table*}[t]\centering\small
\begin{tabular}{llrrrrrrr}\toprule
Workload / trial & Policy & Hot hits & Misses & Loads & Evicts & Register & Seeds & Uncached input\\\midrule
Turnover HTTP/1 & KV & 4279 & 73 & 66 & 66 & 51 & 0 & 15.8\% \\
Turnover HTTP/1 & KV+P & 4551 & 89 & 66 & 66 & 51 & 288 & 8.4\% \\
Turnover HTTP/1 & Domains+P & 4576 & 103 & 88 & 88 & 57 & 327 & 9.7\% \\
Turnover HTTP/1 & Work+P & 4516 & 162 & 144 & 144 & 83 & 326 & 9.3\% \\
Warmed reuse 1 & KV+P & 2954 & 502 & 220 & 218 & 0 & 384 & 13.4\% \\
Warmed reuse 1 & Domains+P & 3203 & 253 & 124 & 124 & 0 & 384 & 14.5\% \\
Warmed reuse 2 & KV+P & 3068 & 388 & 204 & 201 & 0 & 384 & 14.9\% \\
Warmed reuse 2 & Domains+P & 3234 & 222 & 98 & 98 & 0 & 384 & 13.6\% \\
Version turnover 1 & KV+P & 4553 & 87 & 72 & 72 & 51 & 288 & 9.4\% \\
Version turnover 1 & Domains+P & 4566 & 108 & 98 & 98 & 59 & 322 & 9.6\% \\
Version turnover 2 & KV+P & 4524 & 116 & 76 & 76 & 51 & 288 & 9.4\% \\
Version turnover 2 & Domains+P & 4555 & 121 & 101 & 101 & 63 & 324 & 10.0\% \\
Dialogue 1 & KV+P & 4607 & 193 & 112 & 111 & 49 & 192 & 7.6\% \\
Dialogue 1 & Domains+P & 4650 & 150 & 102 & 102 & 41 & 192 & 12.4\% \\
Dialogue 2 & KV+P & 4599 & 201 & 114 & 113 & 49 & 192 & 7.6\% \\
Dialogue 2 & Domains+P & 4694 & 107 & 80 & 80 & 38 & 193 & 11.4\% \\
64 versions 1 & KV & 984 & 1064 & 1168 & 1160 & 0 & 0 & 34.6\% \\
64 versions 1 & Domains & 1864 & 184 & 169 & 161 & 0 & 0 & 45.8\% \\
64 versions 1 & Default & 787 & 1261 & 1806 & 1798 & 0 & 0 & 94.7\% \\
64 versions 2 & KV & 920 & 1128 & 1234 & 1226 & 0 & 0 & 35.1\% \\
64 versions 2 & Domains & 1920 & 128 & 135 & 127 & 0 & 0 & 45.2\% \\
Long 512 & KV+P & 3795 & 68 & 60 & 60 & 39 & 215 & 7.1\% \\
Long 512 & Domains+P & 4321 & 127 & 99 & 99 & 46 & 288 & 10.5\% \\
Long 1,024 & KV+P & 1970 & 5 & 5 & 5 & 1 & 119 & 5.7\% \\
Long 1,024 & Domains+P & 1604 & 34 & 33 & 33 & 2 & 102 & 9.3\% \\
\bottomrule\end{tabular}
\caption{State movement and cache accounting. Hot hits and misses count engine admissions, including auxiliary seeds and retries. Uncached input includes auxiliary seed computation and is divided by logical requested input tokens. Register counts first-use sidecar registration calls, not GPU loads.}\label{tab:state}
\end{table*}

The implementation and experiment launcher are in the Spindle research branch
\texttt{codex/lora-aware-rollout-pool}, attached to
\href{https://github.com/modal-projects/spindle/pull/52}{draft PR 52}.
The base revision before this study is
\texttt{534c3d14e7b856a71ea34af73296c4b3b07319a9}.
The experimental implementation and analysis are recorded in
\texttt{f1a599727dd2515ac6d95191ad2307f09e245694}.
The HTTP/2 connection correction, diagnostics, and settled final observations
are recorded in \texttt{635f53c807378ab4ec546c16eeee7d2e99c05555}.
The routing algorithm itself is unchanged across those revisions.
The first corrected general-serving trial began before the final-observation
settling change. The artifact includes a patch against \texttt{635f53c8} that
reconstructs its exact replay source and verifies the recorded SHA-256.
Revision \texttt{b6b901f} exposes Domains as an explicit sampling-deployment
option; \texttt{d4b2c41} cancels and drains other cohort seed RPCs when one fails,
before releasing local reservations. The latter changes the failure path, not
successful placement or preparation. Per-arm source hashes identify which
helper revision ran.
The deployed SGLang fork is pinned to
\texttt{d050d06437d96196fc68d5b4e5c246408790d537}.

The \texttt{scripts/benchmarks/lora\_routing} directory contains the deployment,
calibration, replay, correctness-control, suite, analysis, and independent
shutdown-monitor programs. \texttt{paper\_suite.py} records every command,
start/end timestamp, and terminal status in a machine-readable manifest.
\texttt{paper\_results.py} validates raw records and emits CSV/JSON metrics and
vector figures. Each final arm records a workload-definition SHA-256 separately
from the hashes of its router, provider, sampler, and replay sources.

The raw artifact directory is \texttt{analysis/pool-paper-20261014}. It contains
successful and failed pilot logs, exact pool manifests, checkpoint provenance,
calibration records, physical-attempt events, engine-counter snapshots, and the
final resource audit. Checkpoint and Modal access are required to repeat the GPU
runs; the summary tables and figures can be rebuilt locally from the saved
records. The source does not embed credentials.

The local analysis environment uses Python 3.12.13, NumPy 2.5.2, Matplotlib
3.11.2, HTTPX 0.28.1, and Modal 1.5.5. Remote images install their declared
dependencies separately; these local versions do not pin every remote package.
The relevant router, sampling, provider, deployment, and cohort-cleanup tests pass
126 tests. Live GPU controls verify preparation on all eight replicas with both
the initial transport and HTTP/2.

An independently running monitor limits this allocation to eight GPUs, below
the requested 48-GPU ceiling. Its initial 180-minute deadline was extended by
45 minutes for the separately controlled HTTP/2 phase, then by 60 minutes to
complete reversed-order trials and longer-output controls after the slow
64-version default run. The resulting bound is 285 minutes and 38 GPU-hours.
The final task-lifetime audit records \textbf{36.71 GPU-hours}, a peak of eight
GPUs, and zero active GPUs after shutdown. GPU task lifetimes span 00:30:12
through 05:06:26 UTC on October 14, 2026. This is a sum of observed task
lifetimes, not a billing invoice. One monitor process exited during the study;
it was restarted with the same deadline and retained task records. That gap and
the deadline amendments are preserved in the artifact rather than described as
uninterrupted monitoring.

The completed controlled HTTP/2 set contains 21 arms and 69,568 logical
requests, including the four longer-output arms. All complete without physical
generation retries, reach at least 288 observed active pool inputs, and pass
the recorded request validation. The main paired conclusions use the eight
order-reversed pairs; longer-output controls and the default-policy arm have
one execution each. The exported latency arrays and rolling curves cover
37 completed arms across all retained phases and pilots.
